% Section 6.8, from lectures/L06/appendix/practice_test/questions.md. The answers are in the
% second block of sv/back/facit/06.tex; the worked solutions are published in
% lectures/L06/appendix/practice_test/solutions.md.

\clearpage
\section{Övningsdugga 1}
\label{c6:sec:ovningsdugga}
Övningsduggan täcker kapitel 1--6 och är en förberedelse inför dugga~1, som är upplagd på samma
sätt och har samma regler. Svaren finns i facit, och fullständiga lösningsförslag i
\file{lectures/L06/appendix/practice_test/solutions.md}.

\paragraph{Tillåtna hjälpmedel} Skrivmaterial och valfri miniräknare, samt den formelsamling som
har använts i kursen. Relevanta formler skrivs upp på tavlan.

\paragraph{Övrigt} Duggan påverkar inte betyget och examinationen i kursen, men den kan ge
0,5~bonuspoäng på den ordinarie tentamen. För att få bonus krävs minst 3,5 poäng av maxpoängen
5,0~poäng. Alla uppgifter kräver lösningar med redovisat svar, inklusive enhet där det behövs.
Lycka till!

\begin{aside}
\note[Obs] Mobiltelefoner får inte användas under den tid som duggan pågår och ska placeras på
angiven plats.
\end{aside}

\prov{od1}

\provuppgift{1}{1,0}
Resistansen för två parallellkopplade resistorer $R_1$ och $R_2$ kan beräknas med formeln
\[
  \frac{1}{R} = \frac{1}{R_1} + \frac{1}{R_2}.
\]
Beräkna resistansen $R_2$ om $R_1 = 10\,\text{k}\Omega$ och parallellresistansen
$R = 4{,}5\,\text{k}\Omega$. Ange svaret i k$\Omega$ med en värdesiffra.

\begin{aside}
\note[Notering] $\Omega$ är enheten för resistans och utläses ”Ohm”.
\end{aside}

\provuppgift{2}{1,0}
Förenkla följande uttryck så långt det går:
\[
  \frac{\dfrac{c^8 - c^4}{c^2}}{c(c^2 - 1)}
\]

\provuppgift{3}{1,0}
Lös nedanstående ekvationssystem. Svara exakt i bråkform.
\[
  \begin{cases}
    4y - 2x + 6 = 0 \\
    3y = 5x + 4
  \end{cases}
\]

\provuppgift{4}{0,5}
Ekvationerna i \hyperref[od1:u:3]{uppgift~3} är exempel på räta linjens ekvation. Ange den första
ekvationens lutning $k$ samt m-värdet $m$.

\provuppgift{5}{0,5}
Ett motstånd med resistansen $R = 4{,}7\,\text{k}\Omega$ genomflyts av strömmen
$I = 2{,}5\,\text{mA}$. Beräkna effekten $P$ med hjälp av tiopotenser:
\[
  P = RI^2
\]
Ange svaret i watt på standardform med två värdesiffror.

\provuppgift{6}{1,0}
Lös följande ekvation:
\[
  8(x + 2)^2 - 5(x - 3)^2 = 4 - 20x
\]
Ange svaret med två decimaler.
